Daily CHIRPS rainfall is added up for each of the 58-day periods.
ERA5-Land has 2\,m air temperature and volumetric soil moisture in the
top layer (0-7\,cm), provided as anomalies relative to a 2015-2021
climatology for the same months. Elevation, slope, and aspect are derived from NASADEM at 30\,m. None
of these are in the automatic label. But the numbers are needed in
the evidence card in case there is justification based on drought or
terrain.

\subsection{Preprocessing and Patch Generation}

Patch generation involves drawing a circle of 320 m radius
around each candidate center (which is half of $64 \times 10$\,m),
creating a bounding square and clipping the square against the
study rectangle, so that patches at the edges are a little bit smaller.
GeoTIFF exports are usually $70 \times 65$ pixels in EPSG:4326 (pixel size $0.0000898^\circ$).
Fifteen rasters are saved per patch for a total of 5,490 files. After cloud masking, the fraction of cloud pixels is 0.8--2.3\%
per patch; candidates with a masked fraction higher than 0.15
pixels in either year were marked as uncertain and not sampled. As
for the VLMs, the T0 and T1 RGB chips were downsampled to
$512 \times 512$ pixels.

The minimum distance between the 366 chosen centers is 639 m
(adjacent non-overlapping squares), while the median is 1,430 m.
The data split is a random stratified sampling procedure instead of a
spatial block split, so spatial autocorrelation between training
and test patches is still an issue.

\subsection{Change Evidence and Labelling}

Let $\Delta$ be the patch-mean NDVI change,
$\mathrm{NDVI}_{T1} - \mathrm{NDVI}_{T0}$. Over the 1,200-candidate
pool, the 15th and 85th percentiles of $\Delta$ are $q_{15} = -0.0377$
and $q_{85} = +0.0031$. The automatic label is
\begin{equation}
\mathrm{class} =
\begin{cases}
\text{uncertain} & \text{mask} > 0.15 \text{ in T0 or T1},\\
\text{decline} & \Delta \le q_{15},\\
\text{increase} & \Delta \ge q_{85},\\
\text{stable} & |\Delta| < 0.05,\\
\text{uncertain} & \text{otherwise}.
\end{cases}
\end{equation}
Since both percentiles fall within $\pm 0.05$, all unmasked patches from
$q_{15}$ to $q_{85}$ are classified as stable and the ``uncertain''
category will not apply to this reserve. The entire area has lost some
greenness (mean $\Delta = -0.016$, s.d.\ $0.027$) and hence ``increase''
refers to those patches that have become greener (relative) rather than
a real improvement in greenness; increase patches have $\Delta$ values of
$-0.001$ to $+0.075$ and half of them have changed by less than 0.011.
``Decline'' refers to those patches that have shown the greatest loss in
greenness ($\Delta$ from $-0.034$ to $-0.071$). Only 34 out of 366 patches
have $|\Delta| \ge 0.05$. The categories are relative rankings in this case.

The first 300 patches (100 per class) were drawn from the candidate
pool. A second batch added up to 25 patches per class, keeping new
centres at least 960\,m (1.5 patch widths) from existing ones; only 22
decline candidates passed this filter, giving 372 patches. Six patches
with incomplete exports were removed, leaving 366: 120 decline, 124
stable and 122 increase. The percentiles for the second batch were
recomputed on the remaining pool, so the class ranges overlap slightly
(decline patches reach $\Delta = -0.034$ and stable patches go down to
$-0.036$).

Every single patch comes with an evidence card with information like
indices, changes, radar measures, climate abnormalities, topography and
pipeline description. This pipeline description contains a string of
text \texttt{candidate\_class=}\ldots which is the field used for
labelling. These trend words (e.g.
\texttt{moderate\_decline}) on the evidence card have been obtained
from $512 \times 512$ tensors using thresholds and not from percentile
rule; these trend words contradict the label in 108 out of 366 cases.
The values for NDVI on the evidence card average to 0.92 of Earth
Engine mean values. Otsu thresholding bounding boxes are present in
only three patches and all narrative explanations are just templates,
because no language model based explanation was obtained during the
dataset construction.

A human reference class was assigned to 19 patches (13 training, 3
development, 3 test). It agrees with the automatic label on 8 of 19
(42\%; Wilson 95\% interval 23--64\%). An earlier
36-patch pilot agreed on 17 of 36 (47\%). There was no second
annotator. The automatic labels are therefore a reproducible ranking,
not a validated disturbance map.

\subsection{Dataset Construction and Data Splits}

Each patch is stored as one JSONL record containing image paths,
per-source statistics, the evidence card and the labels. The patches
were split 256/55/55 into training, development and test sets by
stratified sampling on the label (seed 42). The test set contains 18
decline, 18 stable and 19 increase patches, so always predicting
``increase'' gives a majority-class baseline of $19/55 = 34.5\%$. Every
result in this paper uses this test set.

\subsection{Vision--Language Model Evaluation}

Four instruction-tuned open models were compared:
Qwen2-VL-2B-Instruct ,
LLaVA-Interleave-0.5B ,
PaliGemma-3B-mix-224  and
BLIP-2 OPT-2.7B , all loaded from
Hugging Face. Parameter counts are those of the released checkpoints.
Decoding is greedy. A deterministic parser maps the output to a class:
it reads a single-word answer directly and otherwise searches the text
for the first class word or synonym (for example, ``decrease'' is read
as decline). All 55 outputs were parsed for every model.

\subsubsection{Zero-Shot Evaluation}

The main classification setting, Condition~C, gives a model both RGB
chips and the evidence card and asks for one word: \texttt{decline},
\texttt{increase} or \texttt{stable}. The four models did not receive
exactly the same input (Table~\ref{tab:inputs}). Qwen2-VL-2B and
LLaVA-Interleave-0.5B got images together with a full card, which includes
the \texttt{candidate\_class} summary line. PaliGemma would not respond
to the long prompt for analysis, hence it got images, NDVI summary line,
and a short prompt. The BLIP-2 model takes one image at a time, and hence
the two chips were put together in an image with a short prompt. Zero-shot
tests were run on a laptop.

\begin{table}[t]
\centering
\small
\begin{tabular}{lll}
\toprule
Model & Images & Text \\
\midrule
Qwen2-VL-2B & T0, T1 & full card \\
LLaVA-Interleave-0.5B & T0, T1 & full card \\
PaliGemma-3B & T0, T1 & NDVI line only \\
BLIP-2 OPT-2.7B & stitched & none \\
\bottomrule
\end{tabular}
\caption{Input per model in Condition~C. ``Full card'' includes the
pipeline summary line containing \texttt{candidate\_class}.}
\label{tab:inputs}
\end{table}

Three supporting analyses help interpret the results. First, an
\texttt{evidence\_only} variant removes the summary line and any other
line that names the label. This changes the input only for LLaVA,
because the PaliGemma and BLIP-2 prompts never contained the line, and
Qwen2-VL-2B was run in this mode on only three patches. Second, a
closed-set probe scores LLaVA by the likelihood it assigns to each class
word, with a stated decision rule and a class bias fitted on the
development set. Third, two baselines that use no model apply a
threshold to the $\Delta$NDVI value on the card: one with fixed cut-offs
of $\pm 0.02$, and one with cut-offs fitted on the training set.

\subsubsection{Parameter-Efficient Fine-Tuning}

LoRA  was applied to Qwen2-VL-2B with rank 8,
$\alpha = 16$ and dropout 0.1 on the attention and MLP projection
layers, using AdamW  with a learning rate of
$5 \times 10^{-5}$, weight decay 0.05 and gradient accumulation over
four steps. Training ran on a GTX 1080 Ti node of the institute's Ada
cluster. Two separate runs are reported.
